% Give a critical review of the developments/progress and the current state of the art in the field related to your
% research topic. [ 8-12 pages including visuals) 


% the
\documentclass[../thesis.tex]{subfiles}
\begin{document}

\section{Introduction}
    % talk about security pitfalls of monolithic systems and hybrid systems.
    % AI can expose vulnerabilities
    The only way to have true security and correctness is through formal verification. This fact is of greater importance now than ever due to the prevalence of AI and its ability to find and potentially exploit security vulnerabilities in software. Claude's Mythos model preview has reportedly already found thousands of security vulnerabilities in every major operating system and web browser. \TODO{The number of bugs scales proportionally to the \gls{sloc} quote paper here} and given the scale of the codebases for such software, it is inevitable that as AI becomes more and more capable, the more vulnerabilities will be found. \TODO{insert gernots cycle diagram} The cycle of having a vulnerability exposed then developers rushing to fix it every time is not a real solution to the problem as the delay between fixes will be costly and given the size of the codebases of major software, it is unfeasible to expect that we will eventually fix all the bugs.

    It is especially imperative that operating system and in particular kernel software is correct and secure. This is as even if user applications were to be correct, they inevitably depend on the operating system, and given a faulty operating system, the behaviour of said user application can not be guaranteed. Unfortunately, the major operating systems currently all have massive codebases and therefore many known security vulnerabilities and bugs which makes them unfit for use in security sensitive applications.
    \TODO{talk about the paper: Quantitative Vulnerability Assessment of Systems Software}
    \TODO{Decide whether it would be a good idea to do a deep dive into windows and apple and linux here.}

\section{Operating System}
\begin{figure}[H]
    \centering
    \subfile{diagrams/operating_system.tex}
    \caption{Operating System in a System}
\end{figure}
The operating system is the lowest level of software directly interfacing with hardware and providing user applications a general, usable interface to operate upon. Essentially, it is an abstract machine and a resource manager working with \texttt{policies} to multiplex and allocate hardware resources through a simplified system call interface. \citep{Tanenbaum_Bos:mos5.pdf}. 

\subsection{Kernel}
A key component of any operating system is the operating system kernel. Although kernels are often referenced synonymously with operating systems, there is an important distinction to be made. Whilst parts of the operating system may run on less privileged modes, the kernel is specifically the part of the operating system which runs on the highest level of privilege. This means that the kernel has complete control and access over all computer resources in the system. Regarding the kernel, there are several designs, each having their advantages and disadvantages.
\begin{figure}[H]
    \centering
    \includegraphics{monolithic.png}
    \caption{Types of Kernels}
    \TODO{do stuff here a bout this img. use the svg somehow}
    \TODO{credit: By Golftheman - http://en.wikipedia.org/wiki/Image:OS-structure.svg, Public Domain, https://commons.wikimedia.org/w/index.php?curid=4397379}
\end{figure}

\section{Monolithic Kernel}
A monolithic kernel is a kernel design which has most if not all the operating system components inside the kernel. This usually means that all of these components run in one large address space and the kernel presents a unified system call interface to applications to access hardware resources without having to know specific details. A key advantage of a monolithic kernel is that since everything is inside the kernel, a system call only requires a mode switch into kernel mode unlike the equivalent procedure in a microkernel which requires a context switch. Linux is an example of a monolithic kernel, with all operating system components running in kernel mode. Monolithic kernels like Linux have adopted design elements from microkernels such as loadable kernel modules, which are modules that can be loaded for the operating system during runtime. It is important however to note the key difference being that loadable kernel modules run in kernel mode instead of user mode as would be the case on microkernel systems.

\section{Microkernel}
The converse of the monolithic kernel design is a microkernel which operates on Liedtke's principle of minimality \textit{`A concept is tolerated inside the $\mu$-kernel only if moving it outside the kernel, i.e., permitting competing implementations, would prevent the implementation of the system's required functionality'} \citep{Liedtke_95}. Essentially if a component is not strictly required to be inside the kernel for the system to function, then it should be implemented outside the kernel, i.e. in userland. 

\TODO{talk about the history of the microkernel}
\TODO{talk about microkernels that are prevalent today.}
A recent development is

\TODO{Probably talk about microkernel design, can quote liedtke on microkernel construction, caps and IPC etc. also talk about disadvantages and chen and bershad}
\TODO{also talk about HM as a general purpose operating system.}
\section{Hybrid Kernel}
\begin{figure}[H]
    \centering
    \includegraphics{Windows_2000_architecture.png}
    \caption{\TODO{} By The original uploader was Grm wnr at English Wikipedia.Later versions were uploaded by Xyzzy n at en.wikipedia. - Transferred from en.wikipedia to Commons., CC BY-SA 3.0, https://commons.wikimedia.org/w/index.php?curid=2242280}
    \label{Windows NT Architecture} \TODO{I need to do the svg file instead ideally.}
    
\end{figure}
Many modern operating systems are designed with a mixture of the above design philosophies. This includes Windows NT which implements certain operating system components in user space including \texttt{environment subsystems}, \texttt{server processes}, \texttt{system support processes} and \texttt{NTDLL.DLL} \TODO{cite something and etc. also do better desc of windows.}. 
and the XNU kernel which is based on the Mach microkernel as the core, including the \gls{ipc} infrastructure of Mach as a key component of the \texttt{kernel environment}, with the BSD layer being part of userland providing the application interface \citep{Apple:Kernel_Prog}. 

\section{Djawula}
Currently, monolithic and hybrid kernels dominate the general purpose operating system space and this is despite the software technologically superior design of microkernels \citep{Liedtke_95}. Even the only exception HongMeng includes many design compromises in order to achieve higher performance, to the point that one may argue that HongMeng is no longer a microkernel but instead a hybrid kernel. Having demonstrated that LionsOS can achieve higher performance than Linux on embedded workflows \citep{Heiser_VCJGNLDZAZWPDB_25}, and with a well-designed microkernel (seL4) as the kernel, is it possible to create a performant operating system?

    \section{Exokernels}
    An orthogonal kernel design is the exokernel. The exokernel instead of providing an abstracted interface of the hardware resources in the kernel which even microkernels do through \gls{ipc} and virtual memory, it securely exports hardware resources to untrusted library operating systems, seperating the protection and management of hardware resources to be handled by the exokernel and applications respectively \citep{Engler_KO_95}. The motivation for exokernels is threefold: High level abstractions are costly. Applications know best what resource management designs are optimal for their use cases. Finally, library operating systems offer more flexibility than monolithic operating systems. 

    A criticism that can be made about traditional high level operating system abstractions are that they are overly general. Applications that do not require a feature must pay for substantial overhead which those features incur. \TODO{Lampson and Sproul [32], Anderson et
al. [4] and Massalin and Pu [36]} There is no single best way to implement resource management as the optimality of resource management depends on how applications use these resources \TODO{provide an example of how bad poor resource management can get}. Even when resource management policy is implemented in applications, the hardware abstractions of most kernels hide important information such as page faults and timer interrupts. Even microkernels face this problem, exemplified in the fact that seL4 hides the hardware dirty bit from user applications. \TODO{quote and factcheck} Abstractions also limit the flexibility of applications and this is evident in infeasibility of interesting ideas such as scheduler activations and multiple protection domains within a single address space. 

Following the idea of tailored policies, the exokernel design allows for increased flexibility compared to monolithic kernels, whilst also allowing a high level of detail when implementing these policies compared to microkernels due to a lack of abstractions regarding physical hardware. As the exokernel interfaces to a library operating system, it is conducive to many different library operating systems which can be chosen for the application which will be running, theoretically allowing for optimal performance for the process without extensive reimplementation or modifications to the kernel itself.

Library operating systems as an idea are also highly flexible. Like any operating system, the can be made as compatible and portable as required. Ambitious projects may even decide to skip the implementation of a library operating system and implement their application directly on a exokernel. Since exokernels are untrusted, the library operating system may even decide to trust applications. Similarly to a multi-server operating system design on microkernels, library operating systems can be more easily modified if they are implemented in a similarly modular way. Only the relevant modules are required to be modified in such a scenario.

\TODO{talk about how the design of exokernels are, also about performance and then feasibility as a general purpose operating system and also verifiability}.


% exokernels and other stuff like that.

    \subsection{Unikernel}
    A unikernel can be thought of 
    \section{History of Microkernels}
    % maybe talk about previous designs like size and clans and chiefs.
    \section{State-of-the-Art Microkernels}
    \section{Microkernels for General Operating Systems}

    \TODO{talk about how exokernels could be a similar solution with some tradeoffs}
\end{document}
